% !TeX root = ../main.tex
% Final edited solution based on the revised Markdown and review.tex.
\exercise{۲.۲}
\subsection*{صورت مسئله و فرض‌ها}
نشان دهید در مدل \lr{Shockley}، جریان دو ترانزیستور یکسان سری، هر یک با عرض \(W\) و طول \(L\)، با جریان یک ترانزیستور با عرض \(W\) و طول \(2L\) برابر است.
خواسته اصلی، اثبات برابری در ناحیه \lr{Linear} با شرایط
\(V_{DD}>V_t\) و \mbox{\(0\le V_{DS}<V_{DD}-V_t\)} است؛ کتاب به برقراری نتیجه در ناحیه \lr{Saturation} نیز اشاره می‌کند
\cite[صفحه چاپی ۹۵، شکل \exnum{۲.۳۲}]{weste}.

در هر دو مدار، سورس پایین به زمین و گیت‌ها به \(V_{DD}\) متصل‌اند؛ ولتاژ درین بالایی نسبت به زمین \(V_{DS}\) است.
ولتاژ گره میانی ساختار سری را \(V_1\) می‌نامیم.
تحرک حامل‌ها \(\mu\)، ظرفیت اکسید بر واحد سطح \(C_{\mathrm{ox}}\) و ولتاژ آستانه \(V_t\) برای ترانزیستورها یکسان‌اند.
از \lr{Body Effect} و \lr{Channel Length Modulation} صرف‌نظر می‌شود؛ بنابراین \(V_t\) ثابت و \(\lambda=0\) است.
مدارهای مقایسه‌شده در شکل \ref{fig:series-transistors} بازترسیم شده‌اند.

\begin{figure}[!htbp]
\centering
\begin{LTR}
\begin{circuitikz}[scale=0.92,transform shape]
 \draw (1.8,1.8) node[nmos] (single) {};
 \draw (single.S) -- (1.8,0.5) node[ground] {};
 \draw (single.D) -- (1.8,3.8) node[ocirc] {} node[above] {$V_{DS}$};
 \draw (single.G) -- (0,1.8) node[ocirc] {} node[left] {$V_{DD}$};
 \node[right] at (2.2,1.8) {$W/(2L)$};
 \draw[->] (2.55,3.5) -- (2.55,2.8) node[midway,right] {$I_{DS1}$};
 \node at (1.6,-0.6) {\rl{(الف) ترانزیستور منفرد}};
 \draw (7.6,1.3) node[nmos] (bottom) {};
 \draw (7.6,3.2) node[nmos] (top) {};
 \draw (bottom.S) -- (7.6,0.5) node[ground] {};
 \draw (bottom.D) -- (top.S);
 \draw (top.D) -- (7.6,4.5) node[ocirc] {} node[above] {$V_{DS}$};
 \draw (bottom.G) -- (5.9,1.3) -- (5.9,3.2) -- (top.G);
 \draw (5.9,2.25) -- (5,2.25) node[ocirc] {} node[left] {$V_{DD}$};
 \draw (7.6,2.25) node[circ] {} -- (8.5,2.25) node[right] {$V_1$};
 \node[right] at (8,1.3) {$W/L$};
 \node[right] at (8,3.2) {$W/L$};
 \draw[->] (8.5,4.4) -- (8.5,3.8) node[midway,right] {$I_{DS2}$};
 \node at (7.3,-0.6) {\rl{(ب) دو ترانزیستور سری}};
\end{circuitikz}
\end{LTR}
\caption{مدارهای مقایسه جریان؛ بازترسیم شکل \exnum{۲.۳۲} کتاب \cite[صفحه چاپی ۹۵]{weste}. در هر دو مدار، بدنه‌ها به زمین متصل‌اند؛ اثر بدنه در این تمرین نادیده گرفته می‌شود.}
\label{fig:series-transistors}
\end{figure}

برای کوتاه شدن روابط، کمیت‌های کمکی را به‌صورت زیر تعریف می‌کنیم:
\begin{equation}
 A=V_{DD}-V_t>0,\qquad D=V_{DS},\qquad x=V_1,
 \qquad \beta=\mu C_{\mathrm{ox}}\frac{W}{L}.
 \label{eq:stack-definitions}
\end{equation}
پارامتر ترانزیستور منفرد \(\beta/2\) است؛ ازاین‌رو، جریان آن در \(0\le D<A\)، مطابق مدل خطی کتاب، برابر است با
\cite[رابطه \exnum{۲.۵}، صفحه چاپی ۶۶]{weste}
\begin{equation}
 I_{DS1}=\frac{\beta}{2}\left(AD-\frac{D^2}{2}\right).
 \label{eq:stack-single}
\end{equation}

\subsection*{اثبات در ناحیه خطی و تعیین ولتاژ گره میانی}
ابتدا فرض می‌کنیم هر دو ترانزیستور سری در ناحیه خطی‌اند؛ پس از حل، سازگاری این فرض بررسی می‌شود.
برای ترانزیستور پایین، ولتاژهای گیت‌ـ‌سورس و درین‌ـ‌سورس به‌ترتیب \(V_{DD}\) و \(x\) هستند.
برای ترانزیستور بالا، این ولتاژها به‌ترتیب \(V_{DD}-x\) و \(D-x\) هستند.
بنابراین جریان‌های پایین و بالا عبارت‌اند از
\begin{align}
 I_{\mathrm{b}}&=\beta\left(Ax-\frac{x^2}{2}\right),
 \label{eq:stack-bottom}\\
 I_{\mathrm{t}}&=\beta\left[(A-x)(D-x)-\frac{(D-x)^2}{2}\right]\notag\\
 &=\beta\left(AD-Ax-\frac{D^2}{2}+\frac{x^2}{2}\right).
 \label{eq:stack-top}
\end{align}
در حالت ماندگار، قانون جریان کیرشهف یا \lr{Kirchhoff's Current Law (KCL)} در گره میانی می‌دهد
\(I_{\mathrm{b}}=I_{\mathrm{t}}=I_{DS2}\).
با جمع روابط \eqref{eq:stack-bottom} و \eqref{eq:stack-top}، جمله‌های وابسته به \(x\) حذف می‌شوند:
\begin{equation}
 2I_{DS2}=\beta\left(AD-\frac{D^2}{2}\right)
 \quad\Longrightarrow\quad
 \boxed{I_{DS2}=\frac{\beta}{2}\left(AD-\frac{D^2}{2}\right)=I_{DS1}.}
 \label{eq:stack-equality}
\end{equation}
رابطه \eqref{eq:stack-equality} برابری جریان‌ها را نشان می‌دهد. برای تعیین \(V_1\) طبق راهنمای کتاب، دو جریان را برابر قرار می‌دهیم:
\begin{align}
 Ax-\frac{x^2}{2}&=AD-Ax-\frac{D^2}{2}+\frac{x^2}{2},\notag\\
 x^2-2Ax+AD-\frac{D^2}{2}&=0.
 \label{eq:stack-quadratic}
\end{align}
ریشه‌های معادله \eqref{eq:stack-quadratic} برابرند با
\begin{equation}
 x_{\pm}=A\pm\sqrt{A^2-AD+\frac{D^2}{2}}.
 \label{eq:stack-roots}
\end{equation}
در بایاس مثبت \(0<D<A\)، ترتیب فیزیکی ولتاژ گره‌ها ایجاب می‌کند \(0<x<D\) باشد.
برای بررسی ریشه‌ها، توجه می‌کنیم که
\begin{equation}
 (A-D)^2<A^2-AD+\frac{D^2}{2}<A^2.
 \label{eq:stack-root-bounds}
\end{equation}
هر دو نامساوی از \(D(A-D/2)>0\) نتیجه می‌شوند.
بنابراین \(A-D<\sqrt{A^2-AD+D^2/2}<A\) است.
ریشه مثبت \(x_+>A>D\) پذیرفتنی نیست؛ ریشه منفی در بازه مجاز قرار دارد:
\begin{equation}
 \boxed{V_1=A-\sqrt{A^2-AD+\frac{D^2}{2}},\qquad 0<V_1<D.}
 \label{eq:stack-physical-root}
\end{equation}
برای این ریشه، ترانزیستور پایین شرط \(x<A\) و ترانزیستور بالا شرط \(D-x<A-x\) را دارد.
شرط دوم همان \(D<A\) است؛ همچنین \(A-x>0\) برقرار است.
پس هر دو ترانزیستور روشن و در ناحیه خطی‌اند و فرض آغاز حل تأیید می‌شود.
از خود معادله \eqref{eq:stack-quadratic} نیز داریم
\(Ax-x^2/2=(AD-D^2/2)/2\)؛ جایگذاری آن در رابطه \eqref{eq:stack-bottom}، نتیجه \eqref{eq:stack-equality} را مستقلاً بازتولید می‌کند.

برای کنترل رفتار حدی، بسط ریشه در \(D/A\ll1\) می‌دهد
\begin{equation}
 V_1=\frac{D}{2}-\frac{D^2}{8A}
 +O\!\left(\frac{D^3}{A^2}\right).
 \label{eq:stack-small-voltage}
\end{equation}
پس افت ولتاژ در مرتبه اول تقریباً نصف می‌شود.
در مرز \(D=0\) نیز رابطه \eqref{eq:stack-physical-root} به \(V_1=0\) و هر دو جریان به صفر می‌رسند.

\subsection*{تکمیل نتیجه در ناحیه اشباع}
برای \(D\ge A\)، جریان اشباع ترانزیستور منفرد برابر \(\beta A^2/4\) است.
در ساختار سری، ترانزیستور بالا را در اشباع و ترانزیستور پایین را خطی در نظر می‌گیریم.
برابری جریان‌های این دو ترانزیستور می‌دهد
\begin{equation}
 \beta\left(Ax-\frac{x^2}{2}\right)=\frac{\beta}{2}(A-x)^2
 \quad\Longrightarrow\quad
 x=A\left(1-\frac{1}{\sqrt{2}}\right).
 \label{eq:stack-saturation-root}
\end{equation}
ریشه دیگر از \(A\) بزرگ‌تر است و کنار گذاشته می‌شود.
ریشه انتخاب‌شده شرط خطی بودن پایین، یعنی \(0<x<A\)، و شرط اشباع بالا، یعنی \(D-x\ge A-x\)، را برآورده می‌کند.
جریان شاخه در نتیجه برابر است با
\begin{equation}
 \boxed{I_{DS2}=\frac{\beta}{2}(A-x)^2
 =\frac{\beta A^2}{4}=I_{DS1,\mathrm{sat}}.}
 \label{eq:stack-saturation-current}
\end{equation}
در \(D=A\)، ولتاژ گره و جریان با مقادیر ناحیه خطی پیوسته‌اند.
بنابراین برابری جریان‌ها در هر دو ناحیه و در چارچوب فرض‌های مدل \lr{Shockley} برقرار است.
